Niniejszy rozdział poświęcony jest przede wszystkim wnioskom wyciągniętym z
danych pozyskanych z przeprowadzonego eksperymentu, ale nie ogranicza się
wyłącznie do nich.
Omówione zostały również wnioski wyciągnięte z przeprowadzania samego
eksperymentu.
Idea przyświecająca tej pracy opiera się nie tylko na zbadaniu postawionej
hipotezy, ale również na przyjrzeniu się aspektowi możliwości replikacji
eksperymentu w przyszłości.
Ten aspekt również został omówiony poniżej.

\section{Potwierdzona hipoteza badawcza a praktyka}

Wyniki przedstawione w rozdziale~\ref{cha:wyniki} potwierdzają postawioną
przez praktyków tezę --- a także hipotezę \(H_{1}\) tej pracy ---
w związku z czym istnieje empiryczne potwierdzenie zależności czasu
wykonywania zmian, a przez to i kosztów poniesionych, ze stopniem sprzężenia
klas między sobą.
Wynik eksperymentu pozwala sądzić, że klasy, które mają sztucznie zawyżony
współczynnik \acrshort{MPC} zmodyfikowany poprzez odwróconą refaktoryzację,
są trudniejsze dla programistów do zrozumienia, a tym samym wprowadzane
w nich zmiany obciążone są dodatkowym kosztem.

Sprzężenia opisane przez metrykę \acrshort{MPC} i \acrshort{CBO}
przedstawione w eksperymencie są jednym z jego rodzajów, bardziej
złożone --- takie jak sprzężenie semantyczne --- wymagają dodatkowych badań.

Odpowiedź uzyskana w toku niniejszej pracy, choć eksperymentalnie potwierdzona,
dla praktyków może być jedynie słabą heurystyką, ogólnie potwierdzającą zasadę
,,utrzymywania możliwie najniższego sprzężenia'' podczas implementacji
programu obiektowego.
W praktyce może to oznaczać, że organizacje utrzymujące oprogramowanie
powinny tworzyć oprogramowanie o klasach posiadających najmniejsze możliwe
sprzężenie.
Jednak nie może ona być rozwiązaniem wszystkich problemów związanych
z kosztowną modyfikacją oprogramowania: programy o wielkości milionów
linii kodu mogą wymagać zbyt dużych nakładów pracy, by organizacja
uznała inwestycję w formie refaktoryzacji całego programu za
biznesowo uzasadnioną.
Wszakże istnieją obszary w programach, które nie zmieniają się po
pierwotnym ich wytworzeniu, zatem inwestycja mogłaby nigdy się nie zwrócić.
Wobec tego, organizacje w celu optymalizacji kosztów powinny tworzyć
plany refaktoryzacji poszczególnych obszarów oprogramowania w przyszłości
poddawanych zmianom, jednak takie przewidywanie samo w sobie może okazać
się kosztownym przedsięwzięciem.

Z tego powodu pozostaje otwarte pytanie: jeżeli zakładamy, że istnieje
bazowy (minimalny możliwy) stopień sprzężenia klasy, to jak stwierdzić,
że jest on minimalny?
Obecnie autor pracy nie znalazł odpowiedzi na to pytanie --- prawdopodobnie
jej udzielenie wymaga przeprowadzenia dalszych badań pod tym kątem,
a nawet wprowadzenia teorii minimalnej semantyki implementowanego problemu.

\section{Przeprowadzenie eksperymentu}

Nie wszystkie czynności zaplanowane w tej pracy poszły zgodnie z planem:
niewielka ilość badanych mogła dostarczyć mniej jasnych wyników, a czas
przeprowadzania eksperymentu został niekrytycznie przekroczony w 6 na 10
przypadków.

Dobór programistów był przypadkowy (osoby zgłaszające się do eksperymentu),
a podział na grupy nie uwzględniał

\section{Dalsze badania}

\section{Podsumowanie}